\chapter{Petechiae}

\section*{Definition}
Small (1-3 mm), non-blanching, red or purple spots on the skin caused by intradermal hemorrhage from ruptured capillaries, indicating platelet or vascular abnormalities.

\section*{Pathophysiology}
Petechiae result from bleeding into the dermis due to thrombocytopenia (low platelet count), platelet dysfunction, capillary fragility, or increased intravascular pressure. Non-blanching with diascopy distinguishes petechiae from erythema. Distribution provides clues: dependent areas (lower extremities) suggest vasculitis or increased pressure; widespread suggests systemic thrombocytopenia or coagulopathy.

\section*{Etiology}
\begin{itemize}
  \item Thrombocytopenia (ITP, chemotherapy, leukemia, bone marrow failure)
  \item Platelet dysfunction (aspirin, NSAIDs, uremia)
  \item Vasculitis (Henoch-Schonlein purpura, microscopic polyangiitis)
  \item Infections (meningococcemia, endocarditis, RMSF, viral)
  \item Coagulation disorders (hemophilia, von Willebrand disease)
  \item Trauma or vigorous coughing/vomiting (Valsalva)
  \item Medication reactions
  \item Fat embolism syndrome
  \item Sepsis or disseminated intravascular coagulation (DIC)
  \item Senile purpura (fragile capillaries in elderly)
\end{itemize}

\section*{Indications for Evaluation}
Seek care if:
\begin{itemize}
  \item Severe or worsening symptoms
  \item Petechiae with fever, bleeding from other sites, neurological symptoms, or widespread rapid onset (possible meningococcemia or DIC)
  \item Accompanied by other concerning signs
\end{itemize}

\section*{Related Symptoms}
\begin{itemize}
  \item Easy bruising
  \item Bleeding gums
  \item Nosebleeds
  \item Fever
  \item Fatigue
\end{itemize}

\textbf{Note:} Always consider serious underlying conditions when evaluating persistent petechiae.

\section*{Self-Care / Home Management}
\begin{itemize}
  \item Rest and avoid activities that may aggravate the condition.
  \item Maintain adequate hydration and a balanced diet.
  \item Over-the-counter remedies may provide symptomatic relief (consult a pharmacist or doctor).
  \item Monitor symptoms and seek professional advice if they persist or worsen.
  \item Avoid self-medicating with prescription medications without medical guidance.
\end{itemize}

\section*{Diagnostic Approach}
\begin{itemize}
  \item A thorough history and physical examination are the first steps in evaluation.
  \item Laboratory tests (blood work, urinalysis) may be ordered based on clinical suspicion.
  \item Imaging studies (X-ray, ultrasound, CT, MRI) may be indicated when structural pathology is suspected.
  \item Specialist referral is arranged when the diagnosis remains uncertain or requires specific expertise.
\end{itemize}

\section*{Treatment / Management Overview}
\begin{itemize}
  \item Treatment targets the underlying cause identified through diagnostic evaluation.
  \item Supportive care and symptom management are provided while awaiting definitive therapy.
  \item Medications, lifestyle modifications, and physical therapies may be prescribed as appropriate.
  \item Follow-up ensures treatment effectiveness and allows timely adjustments to the care plan.
\end{itemize}

\clearpage